快速 OT 后端降低了单次 Sinkhorn 更新的成本，但固定宽度批量执行仍会保持完整宽度运行，直到其中最慢的问题完成。我们提出 \policy{}：一个面向批量熵正则最优传输的执行层，它动态移除已经完成的问题，并使后续更新变窄。该方法以独立 Sinkhorn 问题的标准批量执行为起点。一次交替缩放循环后，其中一侧边缘由构造保证为最新；对另一侧边缘进行批量检查，可筛选出需要执行双边际验证器的问题。通过这一已有释放规则的问题会从所有逐问题张量中移除，未完成的问题则继续执行。

分析将计算机会与硬件转化分开。对于首次通过的迭代编号 $n_1,\ldots,n_W$ 和存活宽度 $A_k=|\{t:n_t\geq k\}|$，静态执行消耗 $Wn_{\max}$ 个问题更新槽位，而动态压缩恰好消耗 $\sum_k A_k=\sum_t n_t$。若 $c_k(a,\xi_k)$ 表示宽度为 $a$ 时更新的实测成本，则匹配路径的精确时间差等于宽度成本节省 $\sum_k[c_k(W,\xi_k)-c_k(A_k,\xi_k)]$ 减去增量控制开销。因此，深度异质性创造可移除工作，而硬件宽度成本曲线决定其中多少能够转化为墙钟加速。在实测 A100/Triton 配置上，宽度扫描呈现波次形状的成本曲线：当 $n=1024$ 时，宽度一到六处于一个近似平坦的区域，而更大的问题会更早地响应宽度缩减。

匹配的固定宽度/动态压缩对照在四 worker MetroPT-3 端点上得到 \CMetroCorrectedEndpointSpeedup{}，在 ImageNet-32 训练 solver 字段上得到 $1.054\times$，并在已登记的 MetroPT scale-out 单元中最高达到 $3.110\times$。OTT-JAX 实现在三个容差上得到 $1.366$--$1.581\times$。当多个问题使用同一初始化时，匹配的 packed Triton 动态压缩达到 $1.421\times$，并移除 30.62\% 的问题更新槽位。Packer19 提供匹配的逻辑 mask 对照，以及压缩收益在深度同步时消失的边界。实际残差和已登记应用端点共同构成实现证据；所有报告的消费者质量门和训练质量门均通过。
